\exnumber{07}
\extitle{UI I - Ex 8 - UI Challenge - Watches}
\exfiles{ex07/Figma file, named Name\_FirstName\_DeliverableName\_Date\_VersionNumber}
\exauthorize{Those named in the exercise}
\makeheaderfiles

\textbf{Instructions:}

Copy these images on Figma

Screen1
Screen2
Screen3
Screen4

\textbf{Deliverable:}

\begin{itemize}
\item A Figma file with a copy of the image with Figma elements
\item Make sure:
\begin{itemize}
\item You use symbols and variants
\item All your layers are named and organized
\item You use overlays
\item You use masks to change the shape of the imported image
\end{itemize}
\item Upload the Figma file on Github, labeled as follows: "Name\_FirstName\_DeliverableName\_Date\_VersionNumber". For example, it can be labeled as "Doe\_John\_MidFiPrototype\_05242024\_V1".
\item Add a title within the file.
\item Pay attention to the global aspect of the file. It must be clear, simple and easy to read. You can get inspiration on canva.com but don't overload your design with too much details!
\end{itemize}

\textbf{Tips:}

\begin{itemize}
\item Fonts: SF Compact for the iOS elements,  Futura STD Condensed for the Nike signs
\item The frame for each iWatch screen is watch 44 mm 184 x 224
\item Figma article about Symbols and Variants \url{https://help.figma.com/hc/en-us/articles/360056440594-Create-and-use-variants}
\item You can use existing Libraries by browsing the Figma community resources
\end{itemize}

